\documentclass[11pt,a4paper]{article}

% ---------- Idioma y codificación ----------
\usepackage[spanish,es-tabla]{babel}
\usepackage[utf8]{inputenc}
\usepackage[T1]{fontenc}
\usepackage{lmodern}

% ---------- Formato de página ----------
\usepackage[margin=2.5cm]{geometry}
\usepackage{setspace}
\setstretch{1.15}
\usepackage{parskip}

% ---------- Tablas y listas ----------
\usepackage{booktabs}
\usepackage{tabularx}
\usepackage{array}
\usepackage{enumitem}

% ---------- Títulos en negrita (sin color) ----------
\usepackage{titlesec}
\titleformat{\section}
  {\Large\bfseries}
  {\thesection}{0.8em}{}
\titleformat{\subsection}
  {\large\bfseries}
  {\thesubsection}{0.6em}{}
\titleformat{\subsubsection}
  {\normalsize\bfseries}
  {\thesubsubsection}{0.5em}{}
\titlespacing*{\section}{0pt}{18pt}{8pt}

% ---------- Encabezado y pie ----------
\usepackage{fancyhdr}
\pagestyle{fancy}
\fancyhf{}
\fancyhead[L]{\small\textbf{Scrum --- Sprint \arabic{sprintnum}}}
\fancyhead[R]{\small\textbf{CrediStep}}
\fancyfoot[C]{\small\thepage}
\renewcommand{\headrulewidth}{0.4pt}

% ---------- Número de sprint ----------
\newcounter{sprintnum}
\setcounter{sprintnum}{1}   % Cambie este número en cada sprint

% ---------- Caja del objetivo (solo blanco y negro) ----------
\usepackage[most]{tcolorbox}
\newtcolorbox{objetivo}{
  colback=white, colframe=black,
  boxrule=0.8pt, arc=0pt,
  left=8pt, right=8pt, top=6pt, bottom=6pt,
  fonttitle=\bfseries, title={Objetivo del Sprint}
}

% ======================================================
\begin{document}

\begin{titlepage}
    \centering
    
    {\Large \textbf{Pontificia Universidad Javeriana}} \\
    \vspace{0.2cm}
    {\large Ingeniería de Sistemas} \\
    {\large Gestión de la Innovación} \\
    
    \vspace{3.5cm}
    
    {\Huge \textbf{CrediStep}} \\
    \vspace{0.4cm}
    {\Large Documento Scrum — Sprint 1} \\
    \vspace{0.2cm}
    {\large Contextualización del problema} \\
    
    \vfill
    
    \textbf{Integrantes:} \\
    \vspace{0.2cm}
    Beiker Romario Jaramillo Giraldo \\
    Laura Sofía Ortiz Gómez \\
    Sarahy Michelle Silva Pinto \\
    Valentina Carrillo Peñuela \\
    Miguel Ángel Gómez López \\
    
    \vspace{1cm}
    
    \textbf{Profesor:} Víctor Hugo Vela Pinzón \\
    \vspace{0.8cm}
    Bogotá D.C. \\
    Inicio: 16/07/2026 \\
    Fin: 30/07/2026
    
\end{titlepage}
% ---------- 2. Objetivo del sprint ----------
\section{Objetivo del sprint}

\begin{objetivo}
 Identificar, con fuentes oficiales, el problema que enfrentan los jóvenes de 18 a 25 años para acceder al crédito, y definir el usuario principal al que ayuda CrediStep.
\end{objetivo}

% ---------- 3. Backlog del sprint ----------
\section{Backlog del sprint}

\subsection{Historias de usuario seleccionadas}

\begin{tabularx}{\textwidth}{c X c c c}
  \toprule
  \textbf{ID} & \textbf{Historia de usuario} & \textbf{Prioridad} & \textbf{Puntos} & \textbf{Estado} \\
  \midrule
  HU-01 & Como Product Owner, quiero que el problema esté respaldado con datos oficiales de acceso al crédito de los jóvenes, para confiar en que es real. & Alta & 5 & Terminada \\
  HU-02 & Como Product Owner, quiero el problema y el usuario principal definidos, para saber a quién le ayuda CrediStep. & Alta & 3 & Terminada \\
  HU-03 & Como Product Owner, quiero las fuentes y sus anexos registrados, para poder verificar cada cifra. & Media & 2 & Terminada \\
  \bottomrule
\end{tabularx}

% ---------- 4. Desglose de tareas ----------
\section{Desglose de tareas}

\begin{tabularx}{\textwidth}{c X l c c}
  \toprule
  \textbf{HU} & \textbf{Tarea} & \textbf{Responsable} & \textbf{Horas est.} & \textbf{Estado} \\
  \midrule
  HU-01 & Buscar el Reporte de Inclusión Financiera 2024 & Beiker Jaramillo & 3 & Terminada \\
  HU-01 & Contrastar cifras con BBVA Research y Portafolio & Valentina Carrillo & 4 & Terminada \\
  HU-02 & Redactar el problema y el usuario & Laura Ortiz & 4 & Terminada \\
  HU-02 & Redactar la oportunidad y la propuesta & Sarahy Silva & 3 & Terminada \\
  HU-03 & Organizar fuentes y anexos & Miguel Gómez & 2 & Terminada \\
  \bottomrule
\end{tabularx}

% ---------- 5. Seguimiento diario (Daily Scrum) ----------
\section{Seguimiento diario (Daily Scrum)}

\begin{tabularx}{\textwidth}{c l X X X}
  \toprule
  \textbf{Día} & \textbf{Miembro} & \textbf{Ayer} & \textbf{Hoy} & \textbf{Impedimentos} \\
  \midrule
  1 & Beiker Jaramillo & Revisamos la idea inicial del problema & Buscar los reportes oficiales de inclusión financiera & Ninguno \\
  2 & Valentina Carrillo & Comparamos cifras de crédito con otros estudios & Revisar cada cifra contra el informe original & El informe usa el grupo de 41 a 60 años, no de 41 a 65 \\
  \bottomrule
\end{tabularx}

\subsection{Impedimentos y acciones}

\begin{tabularx}{\textwidth}{X l l}
  \toprule
  \textbf{Impedimento} & \textbf{Responsable} & \textbf{Estado} \\
  \midrule
  Rango de edad mal citado en una cifra & Valentina Carrillo & Resuelto \\
  \bottomrule
\end{tabularx}

% ---------- 6. Métricas del sprint ----------
\section{Métricas del sprint}

\begin{tabularx}{\textwidth}{X c c}
  \toprule
  \textbf{Indicador} & \textbf{Planeado} & \textbf{Real} \\
  \midrule
  Puntos de historia completados & 10 & 10 \\
  Historias completadas          & 3 & 3 \\
  Velocidad del equipo           & 10 & 10 \\
  \bottomrule
\end{tabularx}

% ---------- 7. Revisión del sprint ----------
\section{Revisión del sprint}

\subsection{Incrementos entregados}
\begin{itemize}[leftmargin=1.5em]
  \item {Documento de contextualización con cifras y fuentes verificadas.}
  \item {Descripción del usuario y de la propuesta inicial de CrediStep.}
\end{itemize}

\subsection{Retroalimentación de los interesados}
{Sin comentarios del profesor registrados para este sprint.}

\subsection{Historias no completadas}
\begin{tabularx}{\textwidth}{c X X}
  \toprule
  \textbf{ID} & \textbf{Motivo} & \textbf{Acción propuesta} \\
  \midrule
  -- & Todas las historias se completaron. & No aplica \\
  \bottomrule
\end{tabularx}

% ---------- 8. Retrospectiva ----------
\section{Retrospectiva del sprint}

\subsection{Qué salió bien}
\begin{itemize}[leftmargin=1.5em]
  \item {Encontramos datos oficiales que respaldan el problema: solo el 22 \% de los jóvenes de 18 a 25 años accede a crédito, frente al 35,5 \% general.}
  \item {El problema y el usuario quedaron claros y bien explicados.}
\end{itemize}

\subsection{Qué se puede mejorar}
\begin{itemize}[leftmargin=1.5em]
  \item {La contextualización se documentó después del pitch, cuando debía ser su base.}
  \item {Una cifra se presentó con un rango de edad distinto al del informe (41 a 65 años en lugar de 41 a 60).}
\end{itemize}

\subsection{Acciones de mejora}
\begin{tabularx}{\textwidth}{X l l}
  \toprule
  \textbf{Acción} & \textbf{Responsable} & \textbf{Fecha límite} \\
  \midrule
  Documentar el problema y sus fuentes antes del pitch & Laura Ortiz & 31/08/2026 \\
  Revisar cada cifra contra el informe original & Valentina Carrillo & 31/08/2026 \\
  Guardar los PDF de las fuentes en la carpeta del sprint & Miguel Gómez & 31/08/2026 \\
  \bottomrule
\end{tabularx}

% ---------- 9. Aprobaciones ----------
\section{Aprobaciones}

\vspace{1cm}
\begin{tabularx}{\textwidth}{X X}
  \textbf{Product Owner} & \textbf{Scrum Master} \\[2cm]
  \rule{6cm}{0.4pt} & \rule{6cm}{0.4pt} \\
  Firma y fecha & Firma y fecha \\
\end{tabularx}

\end{document}
